\uuid{7Y3p}
\titre{Équation droite/plan}
\auteur{bodin ; hanani ; mzari}
\organisation{exo7}
\licence{CC BY-NC-SA 4.0}
\datecreate{2020-01-20}
\niveau{L1}
\module{Géométrie}
\chapitre{Géométrie du plan}
\sousChapitre{Distance, norme, produit scalaire}
\difficulte{1}
\theme{équation droite, plan}
\exo7qcmid{lille-1.geometrie-plane.010}

\contenu{
\begin{question}
\qclassification{140.01,140.02}
\qtags{motcle=équation droite/plan}
Soit $D$ la droite passant par les points $A(1,-1)$ et $B(1,0)$. Quelles sont les bonnes réponses ?
\begin{answers}
    \bad{Une équation cartésienne de $D$ est : $x-y+1=0$.}
    \good{Une équation cartésienne de $D$ est : $x-1=0$.}
    \good{$\vec{u}(1,0)$ est un vecteur normal à $D$.}
    \bad{$\vec{u}(1,0)$ est un vecteur directeur de $D$.}
\end{answers}
\begin{explanations}
Les coordonnées de $A$ et $B$ vérifient l'équation $x-1=0$, celle-ci est donc une équation cartésienne de $D$ et $\vec{u}(1,0)$ est un vecteur normal à $D$.
\end{explanations}
\end{question}
}
